\documentclass[11pt]{article}
\usepackage[margin=1in]{geometry}
\usepackage[T1]{fontenc}
\usepackage[utf8]{inputenc}
\usepackage{amsmath,amssymb,amsthm}
\usepackage{enumitem}

\title{ICPC World Finals 2022\\P. Turning Red}
\author{}
\date{}

\begin{document}
\maketitle

\section*{Problem Summary}

Each button press increases the color of every controlled light by $1$ modulo $3$, using the encoding
red $=0$, green $=1$, blue $=2$. Each light is controlled by at most two buttons. We must turn every
light red while minimizing the total number of button presses.

\section*{Key Observations}

\begin{itemize}[leftmargin=*]
    \item Let $x_i \in \{0,1,2\}$ be the number of times button $i$ is pressed, modulo $3$.
    \item For a light controlled by buttons $u$ and $v$, the requirement that it becomes red is one linear
    equation
    \[
    x_u+x_v \equiv \text{need} \pmod 3.
    \]
    If a light is controlled by only one button $u$, we simply get $x_u \equiv \text{need} \pmod 3$.
    \item Therefore each connected component of the graph on buttons, where two buttons are adjacent if
    they control the same 2-button light, can be solved independently.
    \item Inside one component, once the value of one button is fixed, every other button value is forced by
    BFS propagation through the equations.
    \item There are only three possible values for the root button, so we can try all three and keep the valid
    one with minimum total presses.
\end{itemize}

\section*{Algorithm}

\begin{enumerate}[leftmargin=*]
    \item Convert every initial light color into the residue \texttt{need} required to reach red modulo $3$.
    \item Build the incidence structure between buttons and lights.
    \item Traverse the connected components of the button graph.
    \item For each component, try root values $0,1,2$:
    \begin{itemize}[leftmargin=*]
        \item propagate all forced values with BFS;
        \item reject the choice if some equation becomes inconsistent;
        \item otherwise compute the sum of button values in that component.
    \end{itemize}
    \item If no root value works for some component, output \texttt{impossible}. Otherwise sum the best
    component costs.
\end{enumerate}

\section*{Correctness Proof}

We prove that the algorithm returns the correct answer.

\paragraph{Lemma 1.}
Within one connected component of the button graph, fixing the value of one button determines every
other button value uniquely.

\paragraph{Proof.}
Whenever a light is controlled by buttons $u$ and $v$, the equation
\[
x_u+x_v \equiv \text{need} \pmod 3
\]
implies
\[
x_v \equiv \text{need}-x_u \pmod 3.
\]
So knowing one endpoint forces the other. Repeating this along paths determines every button in the
component uniquely. \qed

\paragraph{Lemma 2.}
For a fixed root value in one component, the BFS propagation finds a valid assignment if and only if such
an assignment exists with that root value.

\paragraph{Proof.}
By Lemma 1, every button value is uniquely forced once the root value is fixed. The BFS computes exactly
those forced values. If a contradiction appears, then some light imposes two different values on the same
button, so no valid assignment with this root value exists. If no contradiction appears, all encountered
equations are satisfied, so the computed assignment is valid. \qed

\paragraph{Theorem.}
The algorithm outputs the minimum total number of button presses, or \texttt{impossible} if no solution
exists.

\paragraph{Proof.}
Different connected components share no variables, so the total cost is the sum of independent component
costs. For one component, every feasible solution corresponds to exactly one of the three root values
modulo $3$. By Lemma 2, for each root value the algorithm either proves infeasibility or constructs the
unique valid assignment. Taking the minimum valid cost in each component is therefore optimal, and if
some component has no valid root value then no global solution exists. \qed

\section*{Complexity Analysis}

Each light-button incidence is processed only a constant number of times across all traversals. Therefore
the running time is linear in the input size, namely $O(l+b+\text{total incidences})$, and the memory
usage is the same order.

\section*{Implementation Notes}

\begin{itemize}[leftmargin=*]
    \item A light controlled by zero buttons is immediately fatal unless it is already red.
    \item The optimal number of presses for one button is always one of $0,1,2$, because three extra
    presses have no effect.
\end{itemize}

\end{document}
